\section{Related Work}
\label{sec:related}

\para{Behavioral testability in philosophy of mind.}
The idea that meaningful theoretical distinctions must be behaviorally grounded has deep roots.
\citet{dennett1988quining} argues through a series of thought experiments that qualia---the supposedly private, ineffable properties of experience---are incoherent once stripped of their behavioral and functional correlates.
\citet{turing1950computing} proposed that behavioral indistinguishability should settle questions about machine thought, a position that implicitly endorses the no costume theories rule.
On the other side, \citet{chalmers2018meta} defends the conceivability of zombies---beings behaviorally identical to conscious agents but lacking experience---which, if coherent, would imply that some meaningful distinctions are inherently non-behavioral.
\citet{block1995confusion} draws a related distinction between phenomenal consciousness (P-consciousness) and access consciousness (A-consciousness), arguing that P-consciousness may exist without behavioral access.
Our work does not resolve this philosophical debate but instead asks an empirical question: \emph{to what extent} do current theories actually rely on non-behavioral distinctions?

\para{The unfolding argument.}
\citet{doerig2019unfolding} provide the most direct formal precursor to our work.
They show that any recurrent neural network can be ``unfolded'' into a feedforward network with identical input-output behavior.
Since \IIT assigns high integrated information ($\Phi$) to recurrent architectures but zero to their behaviorally equivalent feedforward counterparts, \IIT's core measure is neither necessary nor sufficient for consciousness as indexed by behavior.
Our Experiment~3 generalizes this argument: rather than proving the existence of a single counterexample, we systematically measure how much each theory's consciousness attributions shift when substrate changes while behavior is held constant.

\para{Theory comparison in consciousness science.}
Several recent efforts have attempted to compare consciousness theories empirically.
The COGITATE adversarial collaboration~\citep{melloni2023adversarial} tested preregistered predictions from \IIT and \GNW across 256 participants using fMRI, M-EEG, and iEEG, but results were inconclusive for both theories.
The ConTraSt database~\citep{yaron2022contrast} cataloged 412 experiments and found that \emph{methodological choices predict theory allegiance irrespective of findings}---a striking pattern consistent with the hypothesis that theories are less empirically distinct than claimed.
\citet{kirkeby2024quantifying} proposed a Bayesian ``Quantification to the Best Explanation'' framework for scoring theories, noting that most theories agree on cognitive predictions but diverge on phenomenal claims.
Our work complements these efforts by directly measuring the behavioral content of each theory's predictions rather than comparing how well each theory fits existing data.

\para{LLMs as proxy reasoners.}
Using language models to systematically apply theoretical frameworks across many scenarios is a relatively new methodological approach.
\citet{butlin2023consciousness} used a theory-heavy approach to derive indicator properties for AI consciousness, consulting theories of consciousness to identify what substrate and processing features might support experience.
We adopt a complementary strategy: rather than asking what each theory requires for consciousness, we ask how consistently each theory's predictions track behavior versus substrate.
The key advantage of an LLM proxy is consistency---it applies each theory with equal engagement, avoiding the bias \citet{yaron2022contrast} documented where human researchers' theory allegiances contaminate their methodological choices.
